At Northern Hemisphere declination of the non-rising stars
$\delta < \varphi - 90^\circ$, where $\varphi$ stands for latitude. So in
Krakow $\delta < -39.9^\circ$. \hfill[2]

Within assumed precession model declination of the star can differ from
ecliptic latitude $\beta$ not more than $\varepsilon = +23.5^\circ$. So the
ecliptic latitude $\beta$ of the stars at the study should be determined.
\hfill[2]

Solving spherical triangle
\[
  \sin\beta = \cos\varepsilon\cdot\sin\delta
    - \sin\varepsilon\cdot\cos\delta\cdot\sin\alpha
\]
\hfill[3]

Substituting by numerical values for the stars at the study one obtains:

Sirius $\beta = -39^\circ\!.7$

Canopus $\beta = -75^\circ\!.9$ \hfill[1]

Due to the precession, declination of Sirius can reach $-63.2^\circ$ so it
will become non rising star. As far as Canopus is concern its declination
never exceed $-52.4^\circ$, so it will never be visible. \hfill[2]
